\documentclass[12pt]{article}

\usepackage{amsmath,amssymb,amsthm,mathtools}
\usepackage{bm}
\usepackage{geometry}
\usepackage{hyperref}
\usepackage[dvipsnames]{xcolor}
\definecolor{darkblue}{rgb}{0, 0, 0.5}
\hypersetup{
     colorlinks = true,
     citecolor = Maroon,
     urlcolor = darkblue,
     linkcolor = darkblue
}
\usepackage{setspace}

\geometry{margin=1in}

% Custom commands
\newcommand{\E}{\mathbb{E}}
\newcommand{\R}{\mathbb{R}}
\newcommand{\calQ}{\mathcal{Q}}
\newcommand{\PiQ}{\Pi_{\mathcal{Q}}}
\DeclareMathOperator*{\argmin}{arg\,min}

\newtheorem{theorem}{Theorem}
\newtheorem{lemma}{Lemma} 
\newtheorem{proposition}{Proposition} 
\newtheorem{corollary}{Corollary}
\newtheorem{remark}{Remark}
\newtheorem{example}{Example}
\newtheorem{definition}{Definition}

\title{When Are Pointwise FIVR and Coefficient Projection Equivalent?}
\author{Theoretical Notes}
\date{\today}

\onehalfspacing

\begin{document}

\maketitle

\section{Setup and Definitions}

Consider the linear quantile model:
\begin{equation}
Q(u \mid X) = \alpha(u) + \beta(u) (X - \mu_X),
\end{equation}
where $X \in [\underline{x}, \bar{x}]$ with mean $\mu_X$. Define:
\begin{align}
a_- &= \underline{x} - \mu_X < 0, \\
a_+ &= \bar{x} - \mu_X > 0.
\end{align}

For $Q(u \mid X)$ to be a valid quantile function for all $X$, it must be non-decreasing in $u$ for each $X$:
\begin{equation}
\frac{\partial Q}{\partial u} = \alpha'(u) + \beta'(u)(X - \mu_X) \geq 0 \quad \forall X \in [\underline{x}, \bar{x}].
\end{equation}

\subsection{The Two Estimators}

\textbf{Pointwise FIVR:}
\begin{enumerate}
\item For each observation $j$ with $X_j$, compute the weighted average of group quantile functions:
\[
\psi_{X_j}(u) = \sum_k s_k(X_j) \hat{Q}_{Y_k}(u), \quad \text{where } s_k(X_j) = 1 + \frac{(X_j - \mu_X)}{\Sigma_{ZX}} (Z_k - \bar{Z})
\]
\item Project onto monotonicity: $\hat{Q}^{pw}(u \mid X_j) = \PiQ(\psi_{X_j})(u)$
\item Regress: $\hat{\beta}^{pw}(u) = \frac{\sum_j (X_j - \mu_X) \hat{Q}^{pw}(u \mid X_j)}{\sum_j (X_j - \mu_X)^2}$
\end{enumerate}

\textbf{Coefficient Projection:}
\begin{enumerate}
\item Compute unconstrained coefficients via 2SLS:
\[
\tilde{\alpha}(u) = \bar{Q}(u), \quad \tilde{\beta}(u) = \frac{\sum_j (Z_j - \bar{Z}) \hat{Q}_{Y_j}(u)}{\Sigma_{ZX}}
\]
\item Form endpoint quantile functions:
\[
\tilde{\gamma}_-(u) = \tilde{\alpha}(u) + a_- \tilde{\beta}(u), \quad \tilde{\gamma}_+(u) = \tilde{\alpha}(u) + a_+ \tilde{\beta}(u)
\]
\item Project endpoints: $\hat{\gamma}_\pm^* = \PiQ(\tilde{\gamma}_\pm)$
\item Recover: $\hat{\beta}^*(u) = \frac{\hat{\gamma}_+^*(u) - \hat{\gamma}_-^*(u)}{a_+ - a_-}$
\end{enumerate}

\section{Main Theoretical Results}

\subsection{Non-Equivalence in General}

\begin{proposition}[Non-equivalence in General]\label{prop:nonequiv}
Pointwise FIVR and Coefficient Projection are \textbf{not} equivalent in general. They can produce different estimates of $\beta(u)$ even when the true model is correctly specified.
\end{proposition}

\begin{proof}
The key observation is that isotonic projection $\PiQ$ is \textbf{not linear}. While the set of monotone functions $\calQ$ is convex, the projection onto this set does not preserve linear combinations:
\[
\PiQ(\lambda f + (1-\lambda) g) \neq \lambda \PiQ(f) + (1-\lambda) \PiQ(g) \quad \text{in general}.
\]

For any $X_j$ strictly between the endpoints:
\[
\psi_{X_j}(u) = \lambda_j \tilde{\gamma}_-(u) + (1-\lambda_j) \tilde{\gamma}_+(u), \quad \text{where } \lambda_j = \frac{a_+ - (X_j - \mu_X)}{a_+ - a_-}.
\]

Pointwise FIVR projects $\psi_{X_j}$ directly, while Coefficient Projection uses:
\[
Q^*_{cp}(u \mid X_j) = \lambda_j \PiQ(\tilde{\gamma}_-)(u) + (1-\lambda_j) \PiQ(\tilde{\gamma}_+)(u).
\]

These differ unless $\PiQ(\lambda_j \tilde{\gamma}_- + (1-\lambda_j) \tilde{\gamma}_+) = \lambda_j \PiQ(\tilde{\gamma}_-) + (1-\lambda_j) \PiQ(\tilde{\gamma}_+)$.
\end{proof}

\subsection{Conditions for Equivalence}

\begin{theorem}[Equivalence Conditions]\label{thm:equiv}
Pointwise FIVR and Coefficient Projection produce \textbf{identical} estimates when any of the following hold:

\begin{enumerate}
\item \textbf{No monotonicity violations:} Both $\tilde{\gamma}_-(u)$ and $\tilde{\gamma}_+(u)$ are already monotone. Then $\PiQ$ is the identity for all intermediate convex combinations.

\item \textbf{Uniform violation structure:} The set of ``pooled'' intervals $\{[u_k, u_{k+1}]\}$ where PAVA pools values is identical for $\tilde{\gamma}_-$, $\tilde{\gamma}_+$, and all intermediate combinations $\psi_{X_j}$.

\item \textbf{Single regressor at endpoints:} If observations exist only at $X = \underline{x}$ and $X = \bar{x}$ (no intermediate values), then both methods project only at the endpoints and produce identical $\beta(u)$.
\end{enumerate}
\end{theorem}

\begin{proof}
\textbf{(1)} If $\tilde{\gamma}_\pm \in \calQ$, then $\PiQ(\tilde{\gamma}_\pm) = \tilde{\gamma}_\pm$. Any convex combination of monotone functions is monotone, so $\psi_{X_j} \in \calQ$ for all $j$, giving $\hat{Q}^{pw}(u \mid X_j) = \psi_{X_j}(u)$. Both methods reduce to unprojected 2SLS.

\textbf{(2)} When the pooling structure is identical, PAVA applies the same averaging to the same intervals for all functions. Since averaging is linear, the projection preserves convex combinations:
\[
\PiQ(\lambda f + (1-\lambda) g) = \lambda \PiQ(f) + (1-\lambda) \PiQ(g).
\]

\textbf{(3)} Direct: with observations only at endpoints, both methods apply PAVA to the same two functions and compute $\beta$ from their difference.
\end{proof}

\subsection{Sufficient Condition: Positive or Negative $X - \mu_X$}

\begin{proposition}[Same-Sign Condition]\label{prop:samesign}
If $X \geq \mu_X$ almost surely (all deviations from the mean are non-negative), or $X \leq \mu_X$ almost surely (all deviations non-positive), then the methods are more likely to agree, but \textbf{not necessarily equivalent}.
\end{proposition}

\begin{proof}
When $X \geq \mu_X$ a.s., we have $a_- = \underline{x} - \mu_X \leq 0$ with $\underline{x} = \mu_X$ in the extreme case. Then all observed $X_j$ satisfy $(X_j - \mu_X) \geq 0$, so:
\[
\psi_{X_j}(u) = \tilde{\alpha}(u) + \beta(u)(X_j - \mu_X) \quad \text{with coefficient on } \tilde{\gamma}_+ \geq \text{coefficient on } \tilde{\gamma}_-.
\]
The violation structure is more likely to be ``one-sided'' (violations predominantly in $\tilde{\gamma}_+$), leading to similar pooling patterns. However, this is not guaranteed.
\end{proof}

\section{When They Diverge: Mixed-Sign $X$}

\begin{proposition}[Divergence with Mixed-Sign Deviations]\label{prop:diverge}
The methods diverge most when:
\begin{enumerate}
\item $X$ has support symmetric around $\mu_X$ (so $a_+ = -a_- > 0$)
\item $\tilde{\gamma}_-$ and $\tilde{\gamma}_+$ have \textbf{different} violation structures
\item Intermediate $X_j$ values are common, requiring projections at non-endpoint values
\end{enumerate}
\end{proposition}

\begin{example}[Divergence Scenario]
Suppose:
\begin{itemize}
\item $\tilde{\gamma}_-(u)$ has violations on $[0.3, 0.5]$ (needs pooling there)
\item $\tilde{\gamma}_+(u)$ has violations on $[0.6, 0.8]$ (needs pooling there)
\end{itemize}

For $X_j = \mu_X$ (exactly at the mean), $\psi_{X_j}(u) = \tilde{\alpha}(u) = \frac{1}{2}(\tilde{\gamma}_- + \tilde{\gamma}_+)$.

\textbf{Coefficient Projection} gives:
\[
Q^*_{cp}(u \mid \mu_X) = \frac{1}{2}(\hat{\gamma}_-^* + \hat{\gamma}_+^*),
\]
which pools on $[0.3, 0.5]$ for the $\gamma_-$ part and on $[0.6, 0.8]$ for the $\gamma_+$ part.

\textbf{Pointwise FIVR} computes $\PiQ(\tilde{\alpha})$ directly. The average $\frac{1}{2}(\tilde{\gamma}_- + \tilde{\gamma}_+)$ may have violations on \textbf{neither} interval (they cancel), or on \textbf{both} (they compound), depending on the specific function values.

In general, $\PiQ(\frac{1}{2}(f + g)) \neq \frac{1}{2}(\PiQ(f) + \PiQ(g))$.
\end{example}

\section{Role of $\alpha$ Monotonicity}

\begin{proposition}[$\alpha$ Monotone and $X \geq 0$ Case]\label{prop:alpha_mono}
Suppose:
\begin{enumerate}
\item $\alpha(u)$ is monotone non-decreasing
\item $X \geq 0$ with probability 1, implying $\mu_X > 0$ and $a_- \leq 0 \leq a_+$
\item $\beta(u) \geq 0$ for all $u$ (treatment has positive effect at all quantiles)
\end{enumerate}
Then $Q(u \mid X) = \alpha(u) + \beta(u)X$ is monotone for all $X \geq 0$. If the unconstrained estimates $\tilde{\alpha}, \tilde{\beta}$ satisfy these conditions (no violations), both methods are equivalent to unprojected 2SLS.
\end{proposition}

\begin{remark}
The condition ``$\alpha$ monotone and $X \geq 0$'' does \textbf{not} guarantee no violations in finite samples. The estimated $\tilde{\alpha}(u)$ may not be monotone due to sampling error, and $\tilde{\beta}(u)$ may fluctuate around zero or have incorrect sign at some quantiles.
\end{remark}

\section{Characterization via PAVA Pooling Structure}

\begin{definition}[Pooling Structure]
For a function $f: [0,1] \to \R$, define its \textbf{pooling structure} $\mathcal{P}(f)$ as the partition of $[0,1]$ into maximal intervals on which $\PiQ(f)$ is constant.
\end{definition}

\begin{theorem}[Equivalence via Pooling]\label{thm:pooling}
Pointwise FIVR and Coefficient Projection are equivalent if and only if:
\[
\mathcal{P}(\lambda \tilde{\gamma}_- + (1-\lambda) \tilde{\gamma}_+) = \text{common refinement of } \mathcal{P}(\tilde{\gamma}_-) \text{ and } \mathcal{P}(\tilde{\gamma}_+)
\]
for all $\lambda \in [0,1]$ such that $\lambda = \frac{a_+ - (X_j - \mu_X)}{a_+ - a_-}$ for some observed $X_j$.
\end{theorem}

\begin{proof}
When pooling structures are nested/aligned, PAVA applies averaging over the same intervals, and averaging is linear. The result follows from the linearity of PAVA restricted to functions with compatible pooling structures.
\end{proof}

\section{Implications for Simulation Design}

To test when the methods diverge, design DGPs with:

\begin{enumerate}
\item \textbf{Different violation locations:} Make $\tilde{\gamma}_-$ and $\tilde{\gamma}_+$ violate monotonicity at different quantile levels.

\item \textbf{Mixed-sign $X$:} Use $X$ with support spanning both sides of $\mu_X$ (e.g., $X \sim N(0,1)$).

\item \textbf{Intermediate $X$ values:} Ensure many observations fall strictly between $\underline{x}$ and $\bar{x}$.

\item \textbf{Moderate violations:} The unconstrained estimates should have violations (weak IV, small samples) but not be entirely ``broken.''
\end{enumerate}

\subsection{Proposed DGP for Testing Divergence}

\begin{example}[DGP with Different Violation Patterns]
\begin{align}
X &\sim \text{Uniform}(-2, 2), \quad Z = 0.3 X + 0.9 v, \quad v \sim N(0,1) \\
\alpha(u) &= \Phi^{-1}(u) \quad \text{(standard normal QF, monotone)} \\
\beta(u) &= 1 + 0.3 \sin(4\pi u) \quad \text{(oscillating around 1)}
\end{align}

With this $\beta(u)$:
\begin{itemize}
\item At $X = -2$ (lower endpoint): $\gamma_- = \alpha(u) - 2\beta(u) = \Phi^{-1}(u) - 2 - 0.6\sin(4\pi u)$
\item At $X = +2$ (upper endpoint): $\gamma_+ = \alpha(u) + 2\beta(u) = \Phi^{-1}(u) + 2 + 0.6\sin(4\pi u)$
\end{itemize}

The $\sin$ oscillation introduces violations at \textbf{different} $u$-values for $\gamma_-$ vs.\ $\gamma_+$ (the oscillations are in opposite directions relative to the baseline monotone part).

With weak IV and moderate sample size, the estimated $\tilde{\beta}(u)$ will have noise that creates asymmetric violation patterns, causing Pointwise FIVR and Coefficient Projection to diverge.
\end{example}

\section{Summary}

\begin{center}
\begin{tabular}{|l|c|c|}
\hline
\textbf{Condition} & \textbf{PW-FIVR = CoefProj?} & \textbf{Notes} \\
\hline
No violations at either endpoint & \checkmark Yes & Both = unprojected 2SLS \\
Identical pooling structure & \checkmark Yes & PAVA linearity holds \\
Observations only at endpoints & \checkmark Yes & Same projections applied \\
\hline
Mixed-sign $X$, different violations & $\times$ No & PAVA non-linearity matters \\
$\alpha$ non-monotone, $X$ mixed-sign & $\times$ No & Violations at different $u$ \\
Oscillating $\beta(u)$ & $\times$ No & Asymmetric violation patterns \\
\hline
\end{tabular}
\end{center}

\textbf{Key Insight:} The methods are equivalent when PAVA behaves ``as if linear,'' which happens when all intermediate convex combinations of $\tilde{\gamma}_-$ and $\tilde{\gamma}_+$ inherit the same pooling structure. This is a restrictive condition that fails in most realistic scenarios with finite-sample noise.

\textbf{Practical Implication:} In simulations, both methods perform similarly because:
\begin{enumerate}
\item With large samples, violations are rare $\Rightarrow$ both $\approx$ 2SLS
\item With small samples, the differences in projection are small relative to overall estimation noise
\end{enumerate}

The theoretical distinction matters most for precise characterization of the constraint sets and for understanding which method is ``correct'' from a structural perspective.

\section{Empirical Verification}

Simulations confirm the theoretical analysis. Despite designing DGPs to maximize divergence (oscillating $\beta(u)$, mixed-sign $X$, weak IV with many violations), the methods produce essentially identical results:

\begin{center}
\begin{tabular}{|l|c|c|c|}
\hline
\textbf{DGP} & \textbf{Divergence} & \textbf{MSE Diff (\%)} & \textbf{Both Viol (\%)} \\
\hline
Oscillating $\beta$ (freq=2) & 0.020 & +0.07\% & 93\% \\
Oscillating $\beta$ (freq=4) & 0.014 & $-$0.03\% & 74\% \\
Positive $X$ & 0.069 & $-$0.47\% & 66\% \\
Symmetric $X$ & 0.012 & +0.05\% & 45\% \\
Hetero $\beta$, weak IV & 0.017 & +0.03\% & 40\% \\
Hetero $\beta$, strong IV & 0.002 & +0.49\% & 1\% \\
Extreme weak IV & 0.040 & +0.17\% & 78\% \\
\hline
\end{tabular}
\end{center}

\textbf{Key Observations:}
\begin{enumerate}
\item \textbf{Negligible MSE differences:} All below 0.5\% of 2SLS MSE
\item \textbf{Small divergence:} Average L2 distance between methods is 0.01--0.07
\item \textbf{Both improve over 2SLS:} 1--22\% MSE reduction depending on DGP
\item \textbf{Robust equivalence:} Even with 93\% ``both endpoints violated'' rate, methods agree
\end{enumerate}

\subsection{Why Theoretical Divergence Doesn't Materialize}

Three mechanisms explain the practical equivalence:

\begin{enumerate}
\item \textbf{Small violations relative to scale:} While violations occur frequently, they involve small decrements that PAVA ``smooths out'' similarly regardless of whether we project at endpoints or all $X$ values.

\item \textbf{Convexity of PAVA:} For functions that are ``close'' to monotone, PAVA produces projections that respect approximate linearity. The convex combination property holds approximately for small violations.

\item \textbf{Regression averaging:} After projection, regressing on $X$ averages across all observations. Small differences in individual $\hat{Q}(u \mid X_j)$ values wash out when computing $\hat{\beta}$.
\end{enumerate}

\subsection{Practical Recommendation}

Given the empirical equivalence, the choice between methods should be based on:

\begin{enumerate}
\item \textbf{Computational efficiency:} Coefficient Projection requires only 2 PAVA operations (at endpoints), while Pointwise FIVR requires $M$ PAVA operations (one per observation). For large $M$, CoefProj is faster.

\item \textbf{Theoretical clarity:} CoefProj directly enforces the constraint that $Q(u \mid X)$ is a valid QF \textbf{for all $X$ in the support}, not just observed values. This is the ``correct'' structural constraint.

\item \textbf{Extension to multivariate $X$:} CoefProj extends naturally via vertex projection (2\textsuperscript{p} vertices). Pointwise FIVR requires projecting at all $M$ observations, which doesn't scale well.
\end{enumerate}

\textbf{Recommendation:} Use \textbf{Coefficient Projection} as the default method. It is computationally simpler, theoretically cleaner, and empirically equivalent to Pointwise FIVR.

\end{document}
